\newpage
\subsection{Operational Semantics}
\subsubsection{Transition Systems}
\inlinedefinition[Transition System] is a tuple $(\Gamma, T, \rightarrow)$, where $\Gamma$ is a set of \bi{configurations},
$T$ is a set of terminal configurations, with $T \subseteq \Gamma$ and $\rightarrow$ is a transition relation, with $\rightarrow \; \subseteq \Gamma \times \Gamma$,
which describes how executions take place. Big-step transitions are of form $\langle s, \sigma \rangle \rightarrow \sigma'$,
e.g. $\langle \texttt{skip}, \sigma \rangle \rightarrow \sigma$

The transition relations are specified as rules of the form ($^*$ optional side-condition, $\varphi_i$ and $\psi$ transitions)
\[
    \begin{prooftree}
        \hypo{\varphi_1}
        \hypo{\dots}
        \hypo{\varphi_n}
        \infer3[(Name)$^*$]{\psi}
    \end{prooftree}
\]
or spelled out, ``If $\varpi_1, \ldots, \varphi_n$ are transitions (and the \textit{side-condition} is true), then $\psi$ is a transition''.

Herein, $\varphi_1, \ldots, \varphi_n$ are called \bi{premises} of the rule and $\psi$ is the \bi{conclusion}. A rule without premises is an \bi{axiom rule}.
